\chapter{数组：优雅的数据存取方式}

\begin{introduction}
  \item 为什么需要数组：一批同类型数据的家
  \item 一维数组的定义、初始化、下标与遍历
  \item 数组在内存中的连续排布
  \item 求最值、逆序、统计、查找等经典算法
  \item 二维数组与矩阵
  \item 数组的边界、退化与常见错误
\end{introduction}

\chaptermotto{「把一百个数据排好队，你只需要记住队首的位置。」}

前面的程序一次只能处理一两个数。可现实中我们经常要处理一批数据：全班 50 个成绩、一帧 128 个采样点、一张 $3 \times 4$ 的矩阵。如果每个数据都用一个变量，代码会变得无法维护。

数组就是为“一批同类型数据”准备的容器。它让数据排成连续的一队，用一个名字和下标来访问。这一章会把数组的“形”和“神”都讲清楚：既会用它，也知道它在内存里长什么样。

\section{为什么需要数组}

假设要读入 10 个成绩并求平均分。不用数组，你可能写成这样：

\begin{codebox}{scores-ugly.c —— 没有数组的笨办法}
#include <stdio.h>

int main(void)
{
    int s1, s2, s3, s4, s5, s6, s7, s8, s9, s10;
    int sum;

    printf("请输入 10 个成绩：");
    scanf("%d%d%d%d%d%d%d%d%d%d",
          &s1, &s2, &s3, &s4, &s5,
          &s6, &s7, &s8, &s9, &s10);

    sum = s1 + s2 + s3 + s4 + s5 + s6 + s7 + s8 + s9 + s10;
    printf("平均分 = %.2f\n", sum / 10.0);
    return 0;
}
\end{codebox}

这段代码能跑，但问题很明显：

\begin{itemize}
  \item 变量太多，读起来累；
  \item 想求最大值、排序、找某个分数，几乎没法写；
  \item 如果成绩变成 100 个，代码就没法维护了。
\end{itemize}

用数组，只需要一个名字 \code{scores} 和一个下标：

\begin{codebox}{scores-array.c —— 用数组读入 10 个成绩}
#include <stdio.h>

int main(void)
{
    int scores[10];
    int i;
    int sum = 0;

    printf("请输入 10 个成绩：");
    for (i = 0; i < 10; i++) {
        scanf("%d", &scores[i]);
    }

    for (i = 0; i < 10; i++) {
        sum = sum + scores[i];
    }

    printf("平均分 = %.2f\n", sum / 10.0);
    return 0;
}
\end{codebox}

\begin{outbox}{运行示例}
请输入 10 个成绩：90 85 77 92 88 60 95 73 81 100
平均分 = 84.10
\end{outbox}

\begin{center}
\fbox{\parbox{0.8\textwidth}{\centering\kaishu
数组 = 一段连续的内存 + 一个名字 + 一组从 0 开始的下标。}}
\end{center}

\section{一维数组}

\subsection{定义}

\begin{center}
\fbox{\parbox{0.62\textwidth}{
\ttfamily 类型 数组名[元素个数];
}}
\end{center}

例如：

\begin{codebox}{define-array.c —— 数组的定义}
int a[10];          /* 10 个 int，下标 0 到 9 */
double b[5];        /* 5 个 double */
char name[20];      /* 20 个 char，第 9 章会用作字符串 */
\end{codebox}

几个要点：

\begin{itemize}
  \item 元素个数必须是\textbf{编译期常量}，不能是变量（C89 不支持变长数组）；
  \item 下标从 \textbf{0} 开始，到 \textbf{元素个数 $-$ 1} 结束；
  \item 数组名是一个整体，代表这段连续内存的首地址（第 8 章会展开）。
\end{itemize}

\begin{c89box}[C89：数组长度必须是常量]
\begin{codebox}{vla.c —— 错误示范（C99 才允许）}
int n = 10;
int a[n];       /* C89 不允许：n 不是编译期常量 */
\end{codebox}
C89 里请写成 \code{int a[10];}。如果确实需要运行时决定长度，应该用 \code{malloc} 动态分配（第 8 章讲指针时介绍），或者定义一个足够大的数组。
\end{c89box}

\subsection{内存布局：连续排布}

数组最重要的性质是：\textbf{所有元素在内存中连续存放}。下图是一个 \code{int a[5]} 的内存示意（假设 \code{int} 占 4 字节，首地址为 0x1000）：

\begin{figure}[htbp]
  \centering
  \begin{tikzpicture}[x=1cm,y=1cm]
    \foreach \i/\v/\addr in {0/10/0x1000,1/20/0x1004,2/30/0x1008,3/40/0x100C,4/50/0x1010} {
      \node[membox,minimum width=1.6cm] (e\i) at (\i*1.8,0) {\v};
      \node[font=\scriptsize,color=black!60,below=1mm of e\i] {a[\i]};
      \node[font=\scriptsize,color=black!45,above=1mm of e\i] {\texttt{\addr}};
    }
    \draw[decorate,decoration={brace,amplitude=4pt},structurecolor!70,line width=0.7pt]
      (-0.8,-1.05) -- (7.4,-1.05);
    \node[font=\small,below=6mm] at (3.3,-1.05) {5 个元素，共 20 字节，连续存放};
  \end{tikzpicture}
  \caption{\texttt{int a[5]} 的内存布局：元素一个挨一个}
  \label{fig:array-memory}
\end{figure}

因为连续，所以第 \code{i} 个元素的地址可以由首地址算出来：

\[
\texttt{\&a[i]} = \texttt{\&a[0]} + i \times \texttt{sizeof(int)}.
\]

这个公式是后面理解指针的关键。现在你只需要记住：\textbf{下标就是“第几个”，地址是“在哪里”，两者可以互相换算}。

\subsection{初始化}

\begin{codebox}{init-array.c —— 数组的初始化}
int a[5] = {10, 20, 30, 40, 50};  /* 全部初始化 */
int b[5] = {1, 2};                /* 部分初始化，其余为 0 */
int c[5] = {0};                   /* 全部为 0 的常用写法 */
int d[]  = {1, 2, 3};             /* 省略长度，编译器数出 3 */
\end{codebox}

\begin{itemize}
  \item 部分初始化时，剩下的元素自动补 0；
  \item \texttt{\{0\}} 是“全部清零”的惯用写法；
  \item 省略长度时，编译器根据初值个数确定大小。
\end{itemize}

\begin{pitfall}[局部数组不会自动清零]
\code{int a[10];} 只是分配了空间，里面的值是随机的“垃圾值”。如果要用，必须先赋值或初始化。只有全局数组和静态数组才会自动清零。
\end{pitfall}

\subsection{下标访问与遍历}

\begin{codebox}{traverse.c —— 数组的遍历}
#include <stdio.h>

int main(void)
{
    int a[5] = {10, 20, 30, 40, 50};
    int i;

    for (i = 0; i < 5; i++) {
        printf("a[%d] = %d\n", i, a[i]);
    }

    return 0;
}
\end{codebox}

\begin{outbox}{运行结果}
a[0] = 10
a[1] = 20
a[2] = 30
a[3] = 40
a[4] = 50
\end{outbox}

遍历的标准写法是：

\begin{codebox}{traverse-template.c —— 遍历模板}
int i;
for (i = 0; i < 5; i++) {
    /* 使用 a[i] */
}
\end{codebox}

注意条件是 \code{i < 5}，不是 \code{i <= 5}。\code{a[5]} 已经越界了！

\subsection{数组大小与元素个数}

\begin{codebox}{array-size.c —— 数组占多少字节}
#include <stdio.h>

int main(void)
{
    int a[10];

    printf("sizeof(a)    = %d\n", (int)sizeof(a));
    printf("sizeof(a[0]) = %d\n", (int)sizeof(a[0]));
    printf("元素个数     = %d\n", (int)(sizeof(a) / sizeof(a[0])));

    return 0;
}
\end{codebox}

\begin{outbox}{运行结果（典型 32 位 int 环境）}
sizeof(a)    = 40
sizeof(a[0]) = 4
元素个数     = 10
\end{outbox}

\code{sizeof(a) / sizeof(a[0])} 是求数组元素个数的经典写法。它只在“数组名还是数组”的时候有效，一旦数组名退化成指针（传给函数后），这个写法就会失效——第 8 章会详细解释。

\section{数组的常见操作}

\subsection{例：求最大值及其下标}

\begin{example}
读入 10 个整数，输出最大值以及它在数组中的下标。
\end{example}
\begin{solution}
思路如下：

\begin{codebox}{max-index.c —— 最大值与下标}
#include <stdio.h>

int main(void)
{
    int a[10];
    int i;
    int max;
    int max_index = 0;

    printf("请输入 10 个整数：");
    for (i = 0; i < 10; i++) {
        scanf("%d", &a[i]);
    }

    max = a[0];
    for (i = 1; i < 10; i++) {
        if (a[i] > max) {
            max = a[i];
            max_index = i;
        }
    }

    printf("最大值 = %d，下标 = %d\n", max, max_index);
    return 0;
}
\end{codebox}

\begin{outbox}{运行示例}
请输入 10 个整数：3 9 1 7 9 5 2 8 6 4
最大值 = 9，下标 = 1
\end{outbox}

如果有多个相同的最大值，这个程序返回第一个。想返回最后一个，把 \code{>} 改成 \code{>=} 即可。
\end{solution}

\subsection{例：数组逆序}

\begin{example}
把数组中的元素原地逆序，例如 \texttt{\{1,2,3,4,5\}} 变成 \texttt{\{5,4,3,2,1\}}。
\end{example}
\begin{solution}
用两个下标，一个从头往后，一个从尾往前，交换后向中间靠拢：

\begin{figure}[htbp]
  \centering
  \begin{tikzpicture}[x=1cm,y=1cm]
    \foreach \i/\v in {0/1,1/2,2/3,3/4,4/5} {
      \node[membox,minimum width=1.4cm] (e\i) at (\i*1.6,0) {\v};
      \node[font=\scriptsize,color=black!55,below=1mm of e\i] {\i};
    }
    \draw[ptr] ([yshift=3mm]e0.north) to[bend left=25] ([yshift=3mm]e4.north);
    \draw[ptr] ([yshift=-3mm]e1.south) to[bend right=25] ([yshift=-3mm]e3.south);
    \node[font=\scriptsize,color=ptrline,above=10mm] at (3.2,0) {交换 a[0] 与 a[4]，a[1] 与 a[3]};
  \end{tikzpicture}
  \caption{双下标向中间靠拢，完成原地逆序}
  \label{fig:reverse-array}
\end{figure}

\begin{codebox}{reverse-array.c —— 原地逆序}
#include <stdio.h>

int main(void)
{
    int a[10];
    int i, j, temp;
    int n = 10;

    printf("请输入 10 个整数：");
    for (i = 0; i < n; i++) {
        scanf("%d", &a[i]);
    }

    for (i = 0, j = n - 1; i < j; i++, j--) {
        temp = a[i];
        a[i] = a[j];
        a[j] = temp;
    }

    printf("逆序后：");
    for (i = 0; i < n; i++) {
        printf("%d ", a[i]);
    }
    printf("\n");

    return 0;
}
\end{codebox}

\begin{outbox}{运行示例}
请输入 10 个整数：1 2 3 4 5 6 7 8 9 10
逆序后：10 9 8 7 6 5 4 3 2 1
\end{outbox}

循环条件 \code{i < j} 保证中间的元素不会被重复交换。
\end{solution}

\subsection{例：统计每个数字出现的次数}

\begin{example}
输入一个非负整数，统计它各个数字（0\textasciitilde9）出现的次数。
\end{example}
\begin{solution}
用数组 \code{count[10]} 当下标为 0\textasciitilde9 的“计数器”：

\begin{codebox}{digit-count.c —— 数字频次统计}
#include <stdio.h>

int main(void)
{
    int n;
    int digit;
    int count[10] = {0};
    int i;

    printf("请输入一个非负整数：");
    scanf("%d", &n);

    if (n == 0) {
        count[0] = 1;
    }
    while (n > 0) {
        digit = n % 10;
        count[digit]++;
        n = n / 10;
    }

    for (i = 0; i < 10; i++) {
        if (count[i] > 0) {
            printf("数字 %d 出现了 %d 次\n", i, count[i]);
        }
    }

    return 0;
}
\end{codebox}

\begin{outbox}{运行示例}
请输入一个非负整数：20260226
数字 0 出现了 2 次
数字 2 出现了 4 次
数字 6 出现了 2 次
\end{outbox}

用“值当下标”是数组最常用的技巧之一。统计、计数、查表都靠它。
\end{solution}

\subsection{例：线性查找}

\begin{example}
在数组中查找某个值，找到则输出下标，找不到输出提示。
\end{example}
\begin{solution}
思路如下：

\begin{codebox}{linear-search.c —— 线性查找}
#include <stdio.h>

int main(void)
{
    int a[10] = {3, 9, 1, 7, 5, 2, 8, 6, 4, 0};
    int key;
    int i;
    int found = -1;

    printf("请输入要查找的数：");
    scanf("%d", &key);

    for (i = 0; i < 10; i++) {
        if (a[i] == key) {
            found = i;
            break;
        }
    }

    if (found >= 0) {
        printf("找到了，下标是 %d\n", found);
    } else {
        printf("没有找到\n");
    }

    return 0;
}
\end{codebox}

线性查找从头到尾依次比较，最坏要比较 $n$ 次。
\end{solution}

\subsection{例：二分查找}

\begin{example}
在一个\textbf{已排序}的数组中查找某个值。
\end{example}
\begin{solution}
二分查找每次比较中间元素，根据大小关系把查找范围缩小一半：

\begin{figure}[htbp]
  \centering
  \begin{tikzpicture}[x=1cm,y=1cm]
    \foreach \i/\v in {0/1,1/3,2/5,3/7,4/9,5/11,6/13,7/15,8/17,9/19} {
      \node[membox,minimum width=1.1cm,minimum height=6mm] (e\i) at (\i*1.25,0) {\v};
      \node[font=\scriptsize,color=black!55,below=1mm of e\i] {\i};
    }
    \node[font=\scriptsize,color=structurecolor,above=3mm of e0] {low};
    \node[font=\scriptsize,color=structurecolor,above=3mm of e9] {high};
    \node[font=\scriptsize,color=ptrline,above=3mm of e4] {mid};
    \draw[ptr] ([yshift=8mm]e4.north) -- ([yshift=4mm]e4.north);
    \draw[decorate,decoration={brace,amplitude=3pt},structurecolor!60]
      ([yshift=-5mm]e0.south west) -- ([yshift=-5mm]e9.south east);
  \end{tikzpicture}
  \caption{二分查找：比较 \texttt{mid}，然后丢掉一半}
  \label{fig:binary-search}
\end{figure}

\begin{codebox}{binary-search.c —— 二分查找}
#include <stdio.h>

int main(void)
{
    int a[10] = {1, 3, 5, 7, 9, 11, 13, 15, 17, 19};
    int key;
    int low = 0;
    int high = 9;
    int mid;
    int found = -1;

    printf("请输入要查找的数：");
    scanf("%d", &key);

    while (low <= high) {
        mid = (low + high) / 2;
        if (a[mid] == key) {
            found = mid;
            break;
        } else if (a[mid] < key) {
            low = mid + 1;
        } else {
            high = mid - 1;
        }
    }

    if (found >= 0) {
        printf("找到了，下标是 %d\n", found);
    } else {
        printf("没有找到\n");
    }

    return 0;
}
\end{codebox}

二分查找的前提是数组\textbf{已经有序}。它的比较次数是 $O(\log n)$，比线性查找快得多：100 万个元素最多比较约 20 次。
\end{solution}

\section{二维数组}

\subsection{定义与初始化}

二维数组可以看成“数组的数组”，例如 \code{int a[3][4]} 表示 3 行 4 列，共 12 个元素。

\begin{codebox}{define-2d.c —— 二维数组的定义与初始化}
int a[3][4];                    /* 3 行 4 列，未初始化 */
int b[2][3] = {                 /* 按行初始化 */
    {1, 2, 3},
    {4, 5, 6}
};
int c[2][3] = {1, 2, 3, 4, 5, 6};   /* 也可以写成一维形式，按行填充 */
int d[3][4] = {{0}};                /* 全部清零 */
\end{codebox}

访问元素用两个下标：\code{a[i][j]}，其中 \code{i} 是行，\code{j} 是列，都从 0 开始。

\subsection{内存布局：按行连续}

二维数组在内存里仍然是\textbf{连续}的，按“先行后列”排列：

\begin{figure}[htbp]
  \centering
  \begin{tikzpicture}[x=1cm,y=1cm]
    \foreach \r/\lab in {0/第 0 行,1/第 1 行,2/第 2 行} {
      \node[font=\small,anchor=east] at (-1.1,-\r*1.0) {\lab};
      \foreach \c in {0,1,2,3} {
        \pgfmathtruncatemacro{\idx}{\r*4+\c}
        \node[membox,minimum width=1.35cm,minimum height=6.5mm]
          (e\r\c) at (\c*1.5,-\r*1.0) {\idx};
      }
    }
    \node[font=\scriptsize,color=black!55] at (2.25,0.8) {线性下标 $0,1,2,\ldots,11$ 按行递增};
  \end{tikzpicture}
  \caption{\texttt{int a[3][4]} 按行存放：第 0 行 4 个，接着第 1 行 4 个……}
  \label{fig:2d-memory}
\end{figure}

元素 \code{a[i][j]} 的地址是：

\[
\texttt{\&a[i][j]} = \texttt{\&a[0][0]} + (i \times 4 + j) \times \texttt{sizeof(int)}.
\]

所以“按行遍历”通常比“按列遍历”更快——因为前者访问的内存是连续的，更符合缓存的工作方式。

\subsection{二维数组的遍历}

\begin{codebox}{traverse-2d.c —— 遍历二维数组}
#include <stdio.h>

int main(void)
{
    int a[3][4] = {
        {1, 2, 3, 4},
        {5, 6, 7, 8},
        {9, 10, 11, 12}
    };
    int i, j;

    for (i = 0; i < 3; i++) {
        for (j = 0; j < 4; j++) {
            printf("%4d", a[i][j]);
        }
        printf("\n");
    }

    return 0;
}
\end{codebox}

\begin{outbox}{运行结果}
   1   2   3   4
   5   6   7   8
   9  10  11  12
\end{outbox}

\subsection{例：矩阵加法}

\begin{example}
输入两个 $3 \times 4$ 的矩阵，输出它们的和。
\end{example}
\begin{solution}
思路如下：

\begin{codebox}{matrix-add.c —— 矩阵加法}
#include <stdio.h>

int main(void)
{
    int a[3][4];
    int b[3][4];
    int c[3][4];
    int i, j;

    printf("请输入矩阵 A（3x4）：\n");
    for (i = 0; i < 3; i++) {
        for (j = 0; j < 4; j++) {
            scanf("%d", &a[i][j]);
        }
    }

    printf("请输入矩阵 B（3x4）：\n");
    for (i = 0; i < 3; i++) {
        for (j = 0; j < 4; j++) {
            scanf("%d", &b[i][j]);
        }
    }

    for (i = 0; i < 3; i++) {
        for (j = 0; j < 4; j++) {
            c[i][j] = a[i][j] + b[i][j];
        }
    }

    printf("A + B =\n");
    for (i = 0; i < 3; i++) {
        for (j = 0; j < 4; j++) {
            printf("%4d", c[i][j]);
        }
        printf("\n");
    }

    return 0;
}
\end{codebox}
\end{solution}

\subsection{例：矩阵转置}

\begin{example}
输入一个 $3 \times 3$ 矩阵，输出它的转置（行列互换）。
\end{example}
\begin{solution}
思路如下：

\begin{codebox}{transpose.c —— 矩阵转置}
#include <stdio.h>

int main(void)
{
    int a[3][3];
    int i, j;

    printf("请输入 3x3 矩阵：\n");
    for (i = 0; i < 3; i++) {
        for (j = 0; j < 3; j++) {
            scanf("%d", &a[i][j]);
        }
    }

    printf("转置后：\n");
    for (i = 0; i < 3; i++) {
        for (j = 0; j < 3; j++) {
            printf("%4d", a[j][i]);
        }
        printf("\n");
    }

    return 0;
}
\end{codebox}

注意输出时用的是 \code{a[j][i]}，而不是 \code{a[i][j]}。
\end{solution}

\section{数组的常见错误}

\begin{center}
\begin{tabular}{@{}p{0.38\textwidth}p{0.52\textwidth}@{}}
\toprule
错误 & 后果 \\
\midrule
下标越界（\code{a[10]} 用于 \code{int a[10]}） & 未定义行为，可能崩溃或悄悄破坏其他数据 \\
数组未初始化就读取 & 得到随机值 \\
用 \code{=} 拷贝数组 & 编译错误；数组不能整体赋值 \\
数组长度用变量 & C89 编译错误 \\
遍历条件写成 \code{i <= n} & 多访问一个元素，越界 \\
\code{sizeof(a)} 用在函数参数上 & 得到的是指针大小，不是数组大小 \\
\bottomrule
\end{tabular}
\end{center}

\begin{pitfall}[数组越界不会报错]
C 语言不会检查数组下标是否越界。\code{a[10]} 会去读写数组后面的内存，可能：
\begin{itemize}
  \item 什么都没发生（最危险，掩盖了 bug）；
  \item 覆盖了别的变量，导致诡异的结果；
  \item 触发段错误，程序崩溃。
\end{itemize}
写循环时一定要反复确认边界：\code{i < n}，不是 \code{i <= n}。
\end{pitfall}

\begin{pitfall}[数组不能整体赋值]
\begin{codebox}{array-assign.c —— 错误：数组整体赋值}
int a[5] = {1, 2, 3, 4, 5};
int b[5];

b = a;      /* 错误：数组名不能作为赋值目标 */
\end{codebox}
正确做法是用循环逐个复制：
\begin{codebox}{array-copy.c —— 逐个复制}
int i;
for (i = 0; i < 5; i++) {
    b[i] = a[i];
}
\end{codebox}
\end{pitfall}

\section{拔高：数组名、退化与栈空间}

\subsection{数组名不是普通的变量}

\code{sizeof(a)} 能得到整个数组的大小，但 \code{a} 在大多数表达式里会“退化”成指向首元素的指针。这就是为什么：

\begin{itemize}
  \item \code{sizeof(a)} 是 40（整个数组）；
  \item 把 \code{a} 传给函数后，函数里的 \code{sizeof(a)} 变成 8（指针大小，64 位环境）。
\end{itemize}

\begin{codebox}{array-decay.c —— 数组名退化}
#include <stdio.h>

void show(int arr[])
{
    printf("函数内 sizeof(arr) = %d\n", (int)sizeof(arr));
}

int main(void)
{
    int a[10];

    printf("main 内 sizeof(a) = %d\n", (int)sizeof(a));
    show(a);
    return 0;
}
\end{codebox}

\begin{outbox}{运行结果（64 位环境）}
main 内 sizeof(a) = 40
函数内 sizeof(arr) = 8
\end{outbox}

所以“在函数里用 \code{sizeof(arr) / sizeof(arr[0])} 求长度”是错的。正确做法是把长度作为参数一起传进去。

\subsection{数组在栈上}

局部数组分配在栈上，大小在编译时确定。栈空间通常只有几 MB，不要定义过大的局部数组：

\begin{codebox}{big-array.c —— 栈溢出的风险}
void f(void)
{
    int a[1000000];   /* 4 MB，可能直接栈溢出 */
    ...
}
\end{codebox}

需要大数组时，可以：
\begin{itemize}
  \item 定义为全局数组或静态数组（放在数据区，不在栈上）；
  \item 用 \code{malloc} 动态分配（第 8 章）。
\end{itemize}

\subsection{变长数组（VLA）}

C99 支持用变量定义数组长度（\code{int a[n];}），但 C89 不支持，而且 VLA 在栈上分配、大小不可控，不建议在嵌段式系统里使用。本书坚持 C89 写法。

\section{综合例题}

\begin{example}
读入 10 个整数，统计其中正数、负数和零的个数。
\end{example}
\begin{solution}
思路如下：

\begin{codebox}{count-sign.c —— 统计正负零}
#include <stdio.h>

int main(void)
{
    int a[10];
    int i;
    int positive = 0;
    int negative = 0;
    int zero = 0;

    printf("请输入 10 个整数：");
    for (i = 0; i < 10; i++) {
        scanf("%d", &a[i]);
    }

    for (i = 0; i < 10; i++) {
        if (a[i] > 0) {
            positive++;
        } else if (a[i] < 0) {
            negative++;
        } else {
            zero++;
        }
    }

    printf("正数 %d 个，负数 %d 个，零 %d 个\n",
           positive, negative, zero);
    return 0;
}
\end{codebox}
\end{solution}

\begin{example}
用数组求斐波那契数列的前 20 项。
\end{example}
\begin{solution}
思路如下：

\begin{codebox}{fib-array.c —— 用数组存斐波那契}
#include <stdio.h>

int main(void)
{
    int f[20];
    int i;

    f[0] = 1;
    f[1] = 1;
    for (i = 2; i < 20; i++) {
        f[i] = f[i - 1] + f[i - 2];
    }

    for (i = 0; i < 20; i++) {
        printf("%d ", f[i]);
    }
    printf("\n");

    return 0;
}
\end{codebox}

用数组的好处是：所有项都保留下来，可以随时回头访问 \code{f[k]}。
\end{solution}

\begin{example}
输入 10 个学生的成绩，输出高于平均分的成绩。
\end{example}
\begin{solution}
思路如下：

\begin{codebox}{above-average.c —— 高于平均分}
#include <stdio.h>

int main(void)
{
    int scores[10];
    int i;
    int sum = 0;
    double average;

    printf("请输入 10 个成绩：");
    for (i = 0; i < 10; i++) {
        scanf("%d", &scores[i]);
        sum = sum + scores[i];
    }

    average = sum / 10.0;
    printf("平均分 = %.2f\n", average);
    printf("高于平均分的成绩：");
    for (i = 0; i < 10; i++) {
        if (scores[i] > average) {
            printf("%d ", scores[i]);
        }
    }
    printf("\n");

    return 0;
}
\end{codebox}

注意 \code{scores[i] > average} 比较的是 \code{int} 和 \code{double}，\code{scores[i]} 会自动转成 \code{double}，没有问题。
\end{solution}

\section{本章小结}

\begin{center}
\begin{tabular}{@{}p{0.26\textwidth}p{0.64\textwidth}@{}}
\toprule
主题 & 要点 \\
\midrule
数组 & 一段连续内存，元素类型相同，下标从 0 开始 \\
定义 & \code{类型 名字[长度];}，长度必须是编译期常量 \\
初始化 & \texttt{\{...\}}；部分初始化补 0；\texttt{\{0\}} 清零 \\
遍历 & \code{for (i = 0; i < n; i++)}，注意边界 \\
内存 & 元素连续存放；下标可换算成地址偏移 \\
二维数组 & \code{a[i][j]}，按行连续存放 \\
常见算法 & 最值、逆序、统计、线性查找、二分查找、矩阵 \\
常见错误 & 越界、未初始化、整体赋值、长度用变量 \\
\bottomrule
\end{tabular}
\end{center}

\begin{tipbox}[下一章预告]
数组只能装同一种类型的数据。如果要描述一个“学生”，需要同时存姓名、学号、成绩，类型各不相同，数组就不够用了。下一章学习\textbf{结构体}——把不同类型的数据打包成一个整体。
\end{tipbox}

\begin{problemset}
  \item 定义 \code{int a[10]}，读入 10 个数并输出，要求每行 5 个。
  \item 读入 10 个数，求最大值、最小值、平均值。
  \item 读入 10 个数，原地逆序后输出。
  \item 读入 10 个数，统计其中偶数和奇数的个数。
  \item 读入一个整数，统计每个数字出现的次数（0\textasciitilde9）。
  \item 在有序数组中用二分查找查找某个数，输出比较次数。
  \item 输入两个 $3 \times 3$ 矩阵，输出它们的和与转置。
  \item 输入 10 个成绩，输出高于平均分的成绩及其下标。
  \item 用数组求斐波那契数列前 30 项（注意是否会溢出）。
  \item 思考题：为什么 \code{sizeof(a) / sizeof(a[0])} 在函数内部求数组长度会失效？应该怎么传长度？
\end{problemset}
